\subsection{Búsqueda Binaria sobre la Respuesta: Maximizar el Mínimo (muralla con refuerzo)}

\graybold{Preguntas guía:}

\begin{itemize}
\item ¿Piden maximizar un mínimo (o minimizar un máximo) tras aplicar UNA modificación elegida de un conjunto de opciones?
\item ¿Puedo verificar rápido, para un valor objetivo fijo, si ALGUNA elección lo logra?
\item ¿La factibilidad es monótona: si un objetivo H es alcanzable, cualquier objetivo menor también lo es?
\end{itemize}

\graybold{Utilidad:} Resolver problemas de "maximizar el mínimo" cuando la modificación disponible tiene una forma estructurada (aquí, un refuerzo en forma de escalera), buscando binariamente el mejor valor objetivo y verificando factibilidad con una condición cerrada en \bigO{n} por candidato.

\graybold{Complejidad:} \bigO{n \log(\max x + K)}.

\graybold{Fundamento teórico:} Para un objetivo $H$ fijo, una posición $i$ con $x_i < H$ tiene una \emph{deficiencia} $d_i = H - x_i$. Si se refuerza en la posición $p$, la posición $i$ recibe bonificación $K-(p-i)$ solo si $p-K+1 \le i \le p$; para que alcance $H$, se necesita $p \le i + K - d_i$. Además, ninguna posición fuera del rango de bonificación puede ser deficiente, lo que obliga a $p \ge R$ (la posición deficiente más a la derecha). Así, el conjunto de posiciones $p$ válidas es la intersección de intervalos $[i,\, i+K-d_i]$ sobre todas las posiciones deficientes $i$; como estos intervalos tienen extremo izquierdo creciente en $i$, la intersección es no vacía si y solo si $R \le K + \min_i(i + x_i) - H$ (tomando el mínimo solo entre las posiciones deficientes). Esta condición se evalúa en \bigO{n}, y como la factibilidad de $H$ es monótona (un refuerzo que logra mínimo $H$ automáticamente logra cualquier $H' < H$), se puede buscar binariamente el mayor $H$ factible.

\graybold{Ejercicio aplicado}

Dada una muralla de N segmentos con alturas iniciales, y un refuerzo que agrega $K$ bloques al segmento elegido, $K-1$ al anterior, y así sucesivamente, hallar la mayor altura mínima posible tras aplicar el refuerzo en un único segmento.

\graybold{I/O:} N ≤ \ten{5} → lectura total con tokenización (patrón 2). Verificado contra fuerza bruta en 300 casos aleatorios pequeños.

\graybold{Código solución}

\begin{lstlisting}[language=Python]
import sys

def solve():
    data = sys.stdin.buffer.read().split()
    n = int(data[0]); k = int(data[1])
    x = list(map(int, data[2:2 + n]))

    def feasible(h):
        # deficientes: posiciones (1-indexadas) por debajo de h
        r = -1          # posicion deficiente mas a la derecha
        best = None      # min(i + x[i]) entre las deficientes (1-indexado)
        for i in range(n):
            if x[i] < h:
                pos = i + 1
                r = pos
                val = pos + x[i]
                if best is None or val < best:
                    best = val
        if r == -1:
            return True          # nadie por debajo de h: cualquier refuerzo sirve
        # interseccion de rangos validos de p: [pos, pos+K-d] para cada deficiente
        # no vacia  <=>  r <= K + best - h
        return r <= k + best - h

    lo, hi = 1, max(x) + k
    while lo < hi:
        mid = (lo + hi + 1) // 2         # busca el mayor H factible
        if feasible(mid):
            lo = mid
        else:
            hi = mid - 1
    print(lo)

solve()
\end{lstlisting}
